Per a obrir i tancar la porta del garatge, disposem d'una cèl·lula fotoelèctrica d'un material alcalí que presenta una funció de treball d'1,20~eV. Sobre la superfície d'aquest material hi fem incidir llum de diverses longituds d'ona: $\lambda_1 = 1,04\ \mu\mathrm{m}$; $\lambda_2 = 0,6\ \mu\mathrm{m}$; $\lambda_3 = 0,5\ \mu\mathrm{m}$.

\begin{apartat}{1,25}
Quina freqüència i quina energia (en eV) tenen els fotons incidents en cada cas?
\end{apartat}

\begin{apartat}{1,25}
Representeu en una gràfica l'energia cinètica màxima dels electrons arrencats del fotocàtode en funció de l'energia dels fotons incidents (en eV). Hi ha electrons arrencats en tots els casos? Justifiqueu la resposta.
\end{apartat}

\dades{%
  $1\ \mathrm{eV} = 1,602\times 10^{-19}\ \mathrm{J}$.\\
  $c = 3,00\times 10^{8}\ \mathrm{m\,s^{-1}}$.\\
  $h = 6,63\times 10^{-34}\ \mathrm{J\,s}$.}
